% [原文 p.16]
\begin{itemize}[leftmargin=1.6em,itemsep=2pt,topsep=2pt]
\item 一般的 6d $(2,0)$ 小弦理论出现在如下极限中：把 IIB 弦论放在与 $\fg$ 相应的离散子群 $\Gamma\subset\mathrm{SU}(2)$ 所对应的 $\C^2/\Gamma$ 奇点上，再取零弦耦合极限 $g_s\to 0$。
\end{itemize}

这些进路告诉我们：圆上的 $\cX(\fg)$ 等价于 5d $\cN=2$ SYM，并且告诉我们在张量分支（tensor branch）\footnote{6d $\cN=(2,0)$ 张量多重态分裂成一个 6d $\cN=(1,0)$ 张量多重态和一个超多重态。张量分支与 Higgs 分支分别是张量多重态或超多重态中的标量场获得 VEV 的真空（在 $(2,0)$ 超对称下两个分支合而为一）。张量分支有时被称为 Coulomb 分支，因为它在 5d 与 4d 中约化为 Coulomb 分支。在 6d $\cN=(1,0)$ 理论中还有矢量多重态，但它们不含标量，故没有相应的分支。}上移动时如何描述 $\cX(\fg)$。

\para{IIA 或 M 理论中的五膜。}我们从 M 理论（等价地，IIA）构造开始，它适用于 $\fg=\fsu(N)$：$\cX(\fsu(N))$ 是 M 理论中一叠 $N$ 个重合 M5 膜的世界体理论，其中已去除解耦的质心自由度。

M 理论是 11 维理论（洛伦兹号差），有 32 个超对称（一个 Majorana 旋量）。它通过种种对偶与更易理解的弦论及超引力相联系，例如它的低能极限由 11 维超引力描述。就我们的目的而言，最有意思的一点是：圆乘以 10 维时空上的 M 理论，等价于该时空上的 IIA 弦。本综述不打算讨论 M 理论与 IIA 及其它弦论之间错综复杂的对偶网，因此我们写得相当示意化。

关于术语的一个标准评注：$p$ 膜是 $(p+1)$ 维客体，有 $p$ 个空间方向、1 个时间方向；例如 M5 膜是 6 维的，具有洛伦兹号差，正合我们的需要。M 理论有两种这样的半 BPS 客体：M5 膜与 M2 膜。同种膜的平坦\footnote{这里我们当作时空是平坦的；膜对几何的反作用不会推翻这些结论。}平行堆叠保持同一半超对称（见习题 2.2），而且在保持膜平坦且平行的前提下移动它们不花费能量。一叠 $N$ 个 D$p$ 膜的世界体理论早已被人们知道是 $p+1$ 维的最大超对称 SYM（见综述 [161]），而 M 理论中膜堆的世界体理论则一直更难确定。

\begin{itemize}[leftmargin=1.6em,itemsep=2pt,topsep=2pt]
\item 一叠重合 M2 膜的世界体理论，现在已知\footnote{视各人的观点而定，本综述中大多数“已知”一词都应换成“猜想”。归根结底，由于物理学家所关心的大多数情形里路径积分尚未被良好定义，几乎所有的非微扰 QFT 结果都是猜测性的。人们可以思考某个结果的“证据”有多少；与许多其它结果相容的结果，就可以作为判断某个数学定义是否可接受的指南。}是 Aharony--Bergman--Jafferis--Maldacena（ABJM）Chern--Simons 物质理论，它是一个具有显式 3d $\cN=2$ 拉格朗日描述的 SCFT，其超对称会增强为膜所保持的预期 3d $\cN=8$ 超共形代数 $\mathfrak{osp}(8|4)$（见综述 [162]）。R 对称性 $\fso(8)$ 转动 M2 膜周围的 11 维空间。它的全息对偶是 $\mathrm{AdS}_4\times S^7$。本综述要说的就这些。

% [原文 p.17]
\item 一叠 $N$ 个重合 M5 膜的世界体理论，就是我们所谓的 $\cX(\fsu(N))$：一个没有拉格朗日描述的 6d $(2,0)$ SCFT。\footnote{除了 $\fg=\mathfrak{a}_{N-1}=\fsu(N)$，还可以通过在 M5 膜上叠加一个 O5 轨形平面来实现 $\fg=\mathfrak{d}_N=\fso(2N)$。}更准确地说，这样给出的是 $\fu(N)$，但膜的 $\fu(1)$ 质心解耦了。R 对称性 $\fso(5)$ 转动 M5 膜周围的空间。全息对偶是 $\mathrm{AdS}_7\times S^4$，它具有预期的对称代数 $\mathfrak{osp}(8^*|4)$，与 3d 情形只差 7d 与 4d 部分中的一些符号。
\end{itemize}

% [Tutor Note: p.16 关于 ABJM 的插话：原文说 ABJM 的 3d $\cN=2$ 超对称"增强为 $\cN=8$"，这对
%  应 M2 膜的情形（Chern--Simons 等级 $k=1$）。严格地说，ABJM 模型 $\mathrm{U}(N)_k\times\mathrm{U}(N)_{-k}$
%  的显式超对称是 $\cN=6$（其 R 对称性是 $\fso(6)=\fsu(4)$），只有在 $k=1,2$ 时才增强到 $\cN=8$
%  （R 对称性 $\fso(8)$，即 $\mathfrak{osp}(8|4)$），$k>2$ 时保持 $\cN=6$。核对方法：数超荷，
%  $\cN=6$ 有 $6\times 4=24$ 个（3d 旋量 2 个实分量，$\cN=6$ 记法下为 12 个复＝24 个实），
%  $\cN=8$ 有 $8\times 4=32$ 个，与 M2 膜的 32 个超对称相符。原文略去了 $k=1,2$ 的限定。]

\para{M 理论构造的推论。}现在考虑圆（乘以五维 Minkowski 空间）上的 $\cX(\fsu(N))$。圆上的 M 理论等价于 IIA 弦论，而绕圆的 M5 膜变成 D4 膜。于是圆上的 $\cX(\fsu(N))$ 等价于 $N$ 个 D4 膜的世界体理论，即 5d $\cN=2$ SYM，正如本节开头所宣布的。量纲分析表明，5d 规范耦合按 $g_{5d}\sim L_5^{1/2}$ 标度，其中 $L_5$ 是紧化圆的长度。

% [Tutor Note: p.17 的 $g_{5d}\sim L_5^{1/2}$ 可独立核对：6d 理论没有无量纲参数，紧化在长度 $L_5$ 的圆上
%  时，5d 规范耦合只能由 $L_5$ 与 6d 的（唯一的）标度组合而成；而 5d 中 $[g_{5d}^2]=$ 长度，
%  故 $g_{5d}^2\propto L_5$，即 $g_{5d}\propto L_5^{1/2}$，与原文一致。这一关系也与 p.12 要点框里的
%  "$R=g_{5d}^2$" 自洽（把符号 $R\leftrightarrow L_5$ 对应起来即可）。]

接着我们由 M 理论构造描述 $\cX(\fg)$ 的真空。超对称真空由 $N$ 个 M5 膜在 5 个横向方向上的位置参数化，并对膜重新标号取商，因为膜是不可区分的。于是真空空间是 $(\R^N)^5/S_N$。在任一一般真空处，所有自由度都有质量（质量正比于膜之间的间距），只有每个单膜附近的涨落例外：已知它们由每张膜一个 6d 阿贝尔 $(B,\lambda,\Phi)$ 理论描述。标量场 $\Phi_I$（$I=6,\dots,10$）描述 $N$ 张 M5 膜各自在横向方向上的涨落（$\fu(1)$ 迹部分除外）。

\para{IIB 弦。}M 理论构造对 $\fg=\mathfrak{a}_{N-1}$ 的 $\cX(\fg)$ 给出很多洞见，也可通过轨形化推广到 $\mathfrak{d}_N$，但它无法实现例外情形 $\mathfrak{e}_6,\mathfrak{e}_7,\mathfrak{e}_8$。为此需要一个对偶的 IIB 描述。

如 [163] 中所理解的那样，对 IIA 理论中一叠 $N$ 个 NS5 膜作横向的 T 对偶，会得到 $A_{N-1}$ 奇点上的 IIB 弦。这一 6d $(2,0)$ 小弦理论的第二种构造（我们不会太多用到它）很容易推广到所有 ADE 情形。把 IIB 弦论放在 Minkowski 空间 $\R^{1,5}$ 乘以由有限子群 $\Gamma\subset\mathrm{SU}(2)$ 得到的商空间 $\C^2/\Gamma$ 上。这样的子群由 ADE 李代数 $\fg$ 分类。例如 $A_{N-1}$ 情形是 $\Gamma=\Z_N$，它在 $\C^2$ 的坐标上作用为 $(z,w)\mapsto(\ee^{2\pi i/N}z,\ee^{-2\pi i/N}w)$。该设置的零耦合极限 $g_s\to 0$ 给出小弦理论，再取零张力极限就给出 6d $(2,0)$ 理论 $\cX(\fg)$。

沿 6d 理论的真空模空间移动，对应于把 $\C^2/\Gamma$ 吹胀成 ALE 空间，即把 $\C^2/\Gamma$ 在原点的奇点消解为一族 $r=\rk\fg$ 个有限大小的二闭链。这些二闭链的大小成为 $\R^{1,5}$ 上 6d 理论的 $r$ 个标量场 $\Phi$，取值于 $\fg$ 的 Cartan 子代数。事实上，它们的 VEV 参数化了 $\cX(\fg)$ 的真空。在任一真空中，红外自由度是：这些

% [原文 p.18]
标量场 $\Phi$、把 IIB 弦论的手征四形式沿每个二闭链积分得到的二形式 $B$，以及一些旋量。我们最终如愿以偿：得到取值为 $\fg$ 的 Cartan 子代数中元素的阿贝尔自对偶二形式规范场多重态 $(B,\lambda,\Phi)$，即 6d $(2,0)$ 理论。

在这一描述中只有 $\mathrm{SO}(4)$ R 对称性显式可见，而且约化到 5d $\cN=2$ SYM 也不容易看出。

% [Tutor Note: p.18 的 orbifold 描述里「只显式看到 $\mathrm{SO}(4)$」的原因：$\C^2/\Gamma$ 的几何只保留
%  $\mathrm{SU}(2)\cong\mathrm{Sp}(1)$ 的局部旋转对称（作为 $\mathrm{SO}(4)=\mathrm{SU}(2)\times\mathrm{SU}(2)$ 的一支），
%  而 6d $(2,0)$ 的充分 R 对称性是 $\fso(5)=\mathfrak{usp}(4)$，故写不出 $\fso(5)$ 的 5 维表示中的 5 个标量。
%  核对方法：$\C^2/\Gamma$ 的等距群最多给出 $\mathrm{SU}(2)$（对 $A_{N-1}$ 还带上 $\Z_N$ 的残余作用），
%  故横向 5 个方向中只有 4 个（即那一支 $\mathrm{SU}(2)$ 的 4 维矢量表示）在超引力近似下可见。]

\subsection{余维 2 与 4 缺陷}

我们回到 $\cX(\fsu(N))$ 的 M 理论构造，考虑在 $\R^{1,10}$ 中沿下列方向放置膜的相交膜构型。

\begin{center}
\begin{tabular}{l*{11}{c}l}
\toprule
M5 & 0 & 1 & 2 & 3 & 4 & 5 & $\cdot$ & $\cdot$ & $\cdot$ & $\cdot$ & $\cdot$ & $\to$ 6d $(2,0)$ 理论 \\
\midrule
M5$'$ & 0 & 1 & 2 & 3 & $\cdot$ & $\cdot$ & 6 & 7 & $\cdot$ & $\cdot$ & $\cdot$ & $\to$ 余维 2 缺陷 \\
M2 & 0 & 1 & $\cdot$ & $\cdot$ & $\cdot$ & $\cdot$ & $\cdot$ & $\cdot$ & $\cdot$ & $\cdot$ & 10 & $\to$ 余维 4 缺陷 \\
\bottomrule
\end{tabular}
\end{center}

每一列是 $\R^{1,10}$ 中的一个方向；一个点表示一叠膜定域在该坐标的某个给定值处，一个数字表示膜沿相应方向伸展。例如 M5 膜位于 $x^6,x^7,x^8,x^9,x^{10}$ 的给定值处，并沿其余所有坐标伸展。加在 M5$'$ 膜上的撇号只用来把后者与 $\cX(\fg)$ 所寄居的 M5 膜区分开。此表中每一叠附加的膜都会破坏一半超对称（见习题 2.2）。同种膜可以有多叠彼此平行，此时它们不再进一步破坏超对称。

M2 膜只沿与 M5 膜横交的一个方向伸展。在这个方向 $x^{10}$ 上，它们可以是无限的、也可以半无限地终止于某张 M5 膜、或者有限地拉伸于两张 M5 膜之间。无论哪种情形，从 $\cX(\fg)$ 的观点看，一叠 M2 膜都插入一个半 BPS 的余维 4 算子，即支撑在 6d 理论中一个二维切片上的算子 [164]。

% [Tutor Note: p.18 的膜表读法（可自查）：M5 的世界体方向是 $012345$，M5$'$ 是 $012367$，二者交叠
%  在 $0123$ 四个方向上，故从 6d 世界体（012345）看，M5$'$ 占据的是 $x^4,x^5$ 之外的 4 个方向，
%  余维为 $6-4=2$；M2 是 $01\,10$，与 M5 只交叠 $x^0,x^1$，故余维为 $6-2=4$。这也解释了正文
%  「M5 与 M5$'$ 看起来在余维 2 相交」的说法：表中的点只表示横向坐标上定域，真实构型是两者
%  融合成光滑复流形（见 Figure 2）。]

这里写出来的样子似乎表明 M5 膜与 M5$'$ 膜在余维 2 相交。实际上它们会融合成一个光滑复流形，后者在大距离上渐近于我们写下的构型。要使这一点发生，当 $x^4,x^5$ 靠近 M5$'$ 膜的位置时，M5 膜的 $x^6,x^7$ 位置应增长到无穷大，如图 2 所示。我们在第 4 节回到具体情形中的讨论。从 $\cX(\fg)$ 的观点看，在大距离处与 M5$'$ 膜的交汇有一个有效描述，即一个 4 维（余维 2）半 BPS 算子 [165]。

当我们在本综述中探究 AGT 对应时，会学到这些缺陷的种种性质，尤其是描述它们的数据。我们发现：

\begin{itemize}[leftmargin=1.6em,itemsep=2pt,topsep=2pt]
\item 余维 2 算子由 $\fg$ 的幂零轨道标记 [7, 100]。在 $\fsu(N)$ 情形，这些相当于 $N$ 的分拆，它指明 $N$ 张 M5 膜在 $x^6,x^7$ 方向上趋于无穷时如何聚成不同的组。额外的连续数据描述这些方向上的长度标度。
\end{itemize}

% [原文 p.19]
\begin{figure}[htbp]
\centering
\includegraphics[width=0.45\textwidth]{fig2.png}
\figtitle{2}{一对沿 $x^4,x^5$ 方向伸展（图中以水平方向表示）的 M5 膜，以及位于 $x^4,x^5$ 平面中某点的一张 M5$'$ 膜的构型。M5 膜与 M5$'$ 膜融合为一个复流形。$x^6,x^7$ 位置（图中以竖直方向表示）在 $x^4,x^5$ 平面中的某一点处发散。图中画的是解耦质心模之后的情形，这也是两族膜对称地发散的原因。}
\end{figure}

\begin{itemize}[leftmargin=1.6em,itemsep=2pt,topsep=2pt]
\item 余维 4 算子由 $\fg$ 的表示标记。对 $\fsu(N)$，回顾每个表示都联系着一张杨图，其中 $\square$ 是基本 $N$ 维表示，$\square\square$ 是对称表示，等等。方块的总数是 M 理论中描述该算子所需的 M2 膜数目。粗略地说，杨图每一行中的方块数表示可以终止在同一张 M5 膜上的 M2 膜数目。
\end{itemize}

% [Tutor Note: p.19 正文里的杨图符号，原文是排版好的实心小方块，本译本用 amssymb 的 \square
%  与并排的两个 \square 复刻「一格＝基本表示、一行两格＝对称表示」。核对方法：su(N) 的基本表示为
%  一格、对称表示为一行两格、反对称表示为竖排两格（原文未画，故正文只出现前两种）；
%  方块总数等于 M2 膜数，它由表示的最高权（行数/行长）决定，与 g 的秩无直接关系。]

\begin{exercise}{2.2}
平坦的 M5 膜沿 $x^0,x^1,\dots,x^5$ 方向时保持满足 $\Gamma^{012345}\epsilon=\epsilon$ 的超对称；而平坦的 M2 膜沿 $x^0,x^1,x^{10}$ 方向时保持满足 $\Gamma^{23456789}\epsilon=\epsilon$ 的超对称。验证上面的膜构型满足：每一叠附加的膜都破坏一半超对称。（提示：验证 $\Gamma^{01}$、$\Gamma^{23}$、$\Gamma^{45}$ 等彼此对易。）膜堆之间还有哪些相对取向能保持一半超对称？
\end{exercise}

% [Tutor Note: p.19 习题 2.2 中 M2 的投影算子为什么写成 $\Gamma^{23456789}$：11 维旋量满足手征条件
%  $\Gamma^{0\,1\cdots10}\epsilon=\pm\epsilon$。把该乘积按指标分成 $\{0,1,10\}$ 与 $\{2,\dots,9\}$ 两组，
%  并利用 M5 的投影 $\Gamma^{012345}\epsilon=\epsilon$，可把 M2 沿 $x^0,x^1,x^{10}$ 的投影
%  $\Gamma^{0\,1\,10}\epsilon=\epsilon$ 等价地改写为 $\Gamma^{2\cdots9}\epsilon=\epsilon$（两者只差一个整体符号）。
%  核对方法：$\Gamma^{0\cdots10}=\pm\Gamma^{0\,1\,10}\Gamma^{2\cdots9}$，两边作用在 $\epsilon$ 上并代入
%  $\Gamma^{0\,1\,10}\epsilon=\epsilon$ 即得。这正是原文提示中"$\Gamma^{01},\Gamma^{23},\Gamma^{45}$ 彼此对易"的由来：
%  每个两指标因子平方为 $\pm1$，故每叠附加膜都把超荷数减去一半。]

\section{由 6d 得到的 Class S 理论}

\begin{keybox}{从 6d 看 4d $\cN=2$ class S 理论}
\begin{itemize}[leftmargin=1.4em,itemsep=1pt,topsep=1pt]
\item 把 6d 理论 $\cX(\fg)$ 作一个扭曲后约化在带穿刺点的黎曼面 $C$ 上，就得到 4d $\cN=2$ 理论。
\item 它的 Coulomb 分支由微分 $\phi_k=u_k\dd z^k$ 参数化，其阶数 $k$ 与 $\fg$ 的 Casimir 不变量的阶数相同。
\item 它的 SW 曲线 $\Sigma\subset T^*C$ 是 $C$ 的多重覆盖。
\item 粘合穿刺点相当于把对称性规范化。
\item 切开 $C$ 相当于让规范场解耦。
\end{itemize}
\end{keybox}

我们接下来的任务是把 6d 理论 $\cX(\fg)$ 在黎曼面 $C_2$ 上作维数约化。我们在小节 3.1 解释一个部分拓扑扭曲，使得约化后的理论具有 4d $\cN=2$ 超对称 [1]。给出红外物理的 Coulomb 分支与 SW 曲线将在小节 3.2 与小节 3.3 中给出。然后我们在小节 3.4 解释 4d 理论如何分解为一些叫做 tinkertoy 的积木 [95]。

% [Tutor Note: p.19 的要点框在原文里同样贴右缘绕排（左栏正文只有约 170pt 宽，见该页正文行宽），
%  译文改为先出要点框、再接正文段落。此处 $\phi_k=u_k\dd z^k$ 的 $\dd z^k$ 中 $k$ 是上标（$\dd z$ 的
%  $k$ 次幂，即 $k$ 阶微分），不是下标；原文 p.23 的 $\mathcal{O}_j$、$P_k(\Phi_z)\dd z^{d_k}$ 用的是
%  同一约定，可交叉核对。tinkertoy 一词保留英文（见术语表）。]
